\section{Introduction}\label{introduction}

The case for reorganizing introductory statistics around
simulation-based reasoning rather than closed-form normal-theory
derivation is, by now, well established. \citet{cobb2007}'s
analysis of the introductory curriculum, the
simulation-based-inference movement organized around Tintle and
colleagues \citep{tintle2014}, and the GAISE College Report
\citep{gaise2016} all argue that students develop more durable
inferential intuitions when they construct a sampling distribution
empirically before they meet it as a formula. Empirical evaluations
support this conclusion \citep{chance2004, delmas2007}. The
pedagogical principle is now mainstream; what remains uneven is the
\emph{toolkit} through which it is delivered.

The current landscape of simulation-based teaching tools is
fragmented: \emph{Sampling SIM} and \emph{Tinkerplots} address
sampling-distribution intuition \citep{konold2008}; \emph{Tinkerplots}
and \emph{CODAP} address school-level exploratory analysis
\citep{finzer2013}; \emph{NetLogo} and the
\emph{seeing-theory.brown.edu} / \emph{RossmanChance.com} libraries
address agent-based modeling and specific inferential techniques
respectively \citep{wilensky1999}; R Shiny supports
instructor-built applets at considerable authoring cost. Each tool
is highly effective in its specific domain, but each is also
\emph{siloed}: a self-contained applet whose output cannot be piped
into the next analytical step. This matters because the inferential
workflow students need to internalize ---
\emph{data-generating process $\rightarrow$ sampling $\rightarrow$
estimation $\rightarrow$ inference $\rightarrow$ diagnostics} ---
is exactly the kind of composable pipeline that
one-applet-per-concept tools render invisible.

We describe \textbf{Empiria}, an open-source teaching suite for VCV
Rack 2 \citep{vcv-rack} --- the free modular-synthesizer environment
--- designed to address this gap. The \textbf{Methods} plugin (the
subject of this paper) places fifteen statistical primitives into
VCV Rack's patch-cable grammar so that the inferential workflow
itself becomes a left-to-right signal chain on the screen. Five
anchor modules --- \texttt{Sample}, \texttt{Frame}, \texttt{Regress},
\texttt{Test}, \texttt{Boot} --- carry the
simulation-based-inference argument; ten supporting modules
(\texttt{Lag}, \texttt{Code}, \texttt{Tab}, \texttt{Strata},
\texttt{Cohort}, \texttt{Factor}, \texttt{Seed}, \texttt{Tape},
\texttt{Gauge}, \texttt{Quantity}) extend the pipeline to
autocorrelation diagnostics, contingency-table inference,
decomposition, online PCA, deterministic seeding, record-and-replay,
and unit interpretation. Every module ships with a live
visualization matched to the statistic it computes --- a histogram
with a confidence band, a scatterplot with a fitted line, a null
distribution with shaded rejection regions, a bootstrap distribution
annotated with bias-correction diagnostics --- so the dynamic being
taught is visible as it unfolds.

This paper makes the pedagogical and methodological case for
Empiria's design and demonstrates the platform through three worked
classroom examples that span the full range of intended audiences
--- from a 45-minute general-audience demonstration that requires no
prior statistical background, to a precisely seeded reproducible
analysis that satisfies graduate-level scrutiny. Per-module
reference, complete parameter listings, numerical-recipe details,
and patch-grammar documentation are provided in the Methods plugin
manual (\texttt{docs/methods\_manual.md}) and the project
repository; this paper concentrates on the conceptual moves and on
the worked examples through which Empiria can be evaluated as a
teaching tool.

\section{Voltage as the medium of simulation}\label{voltage-as-medium}

Empiria's design rests on a particular intuition about voltage. In
a modular synthesizer, every signal --- whether it represents a
pitch, an amplitude, a gate, or, in our case, a population
fraction or an estimator's value --- is a \emph{control voltage}
(CV), a single floating-point number streaming through a patch
cable. This generic ``unit of stuff'' is the platform's universal
medium, and that universality is exactly what makes it useful for
teaching.

The mapping to a teaching context is direct:

\begin{itemize}
\tightlist
\item A voltage \textbf{is} a real-valued state variable.
\item A patch cable \textbf{is} a wire carrying that variable to
      another operator.
\item A knob \textbf{is} a parameter tuning the operator's response.
\item A module \textbf{is} a transformation --- filter, sum,
      threshold, random draw, regression, hypothesis test.
\item A patch \textbf{is} a system of equations, composed by hand
      from these operators, visible end-to-end.
\end{itemize}

In statistics this mapping is immediate. A voltage stream from
\texttt{Sample} \emph{is} a sample path from a parametric
distribution; running it through \texttt{Frame} \emph{is} the act of
measurement; running it through \texttt{Regress} \emph{is} the
construction of an estimator; running it through \texttt{Test}
\emph{is} a hypothesis test. Each stage in the chain is a module,
and the chain itself is a patch cable: the inferential workflow is
no longer abstract, it is laid out left-to-right on the screen.

The convention is not a synth-builder's aesthetic. It is a direct
inheritance from the analog computers of the
1930s--1960s --- machines such as Vannevar Bush's Differential
Analyzer \citep{bush1931} that solved differential equations by
letting voltages represent state variables, with electrical
components implementing the operators that combined them. These
machines were used to model missile trajectories, weather systems,
predator--prey dynamics, and the spread of epidemics
\citep{small2001} --- substantive applications that overlap
substantially with Empiria's pedagogical scope. The argument here
is not historical nostalgia: it is that voltage proved to be a
remarkably general representation of state, and operating on
voltage proved to be a remarkably general way to compute. Digital
simulation of those same continuous voltages --- which is what VCV
Rack does, at audio sample rate --- inherits the same generality at
a fraction of the cost.

Concrete-to-abstract sequencing has also long been argued to
scaffold the learning of mathematical concepts.
\citet{bruner1966} articulated the enactive $\rightarrow$ iconic
$\rightarrow$ symbolic progression at the level of general theory of
instruction; \citet{sfard1991} framed the same trajectory as the
operational-to-structural duality of a mathematical idea;
\citet{mayer2009}'s cognitive theory of multimedia learning provides
the corresponding empirical foundation for animation-plus-narration
over text-only presentation. Empiria's patch grammar makes this
sequence operational: a two-sample \emph{t}-test is, in Empiria, a
literal sequence of objects the student wires together --- two
\texttt{Sample} modules, a \texttt{Frame} on each branch, a
\texttt{Test} consuming both --- and only once that patch has been
built and modulated does the student meet the corresponding formula
on paper, at which point each algebraic symbol corresponds to a
module already manipulated.

\section{Five anchor modules}\label{five-anchor-modules}

The Methods plugin's fifteen modules share a common 20HP panel
grammar: a header strip carrying the module title, a 280 $\times$ 190
px visualization area, a row of knobs for primary parameters, a row
of trimpots or buttons for set-once parameters, and rows of inputs
and outputs. We describe the five \textbf{anchor modules} in
narrative form; full per-module reference for these and the ten
supporting modules is in the plugin manual.

\textbf{\texttt{Sample}} is the entry point of every Methods
pipeline. It draws an i.i.d.\ stream from a user-selected parametric
distribution --- Normal, Uniform, Exponential, or Beta --- at a
user-supplied clock rate, and emits the running sample mean and
standard deviation alongside the raw draw. The live panel overlays
the theoretical PDF on the empirical histogram, so a student
literally watches the empirical shape converge to the theoretical
as \emph{n} grows. Right-click presets re-label the same draws as
substantive quantities --- adult human height, IQ score, reaction
times, right-skewed income, U-shaped opinion --- making the
parametric family pedagogically useful: the same mathematical
object becomes height, IQ, or income with one menu click.

\textbf{\texttt{Frame}} is the bridge between the i.i.d.\ stream and
inferential statistics. It collects samples on a clock into a
buffer of configurable length and reports the sample mean, sample
SD, and the standard error of the mean. Three modes are exposed via
the front-panel switch: \textbf{SNAPSHOT} collects \emph{n} and
then freezes (a cross-sectional view); \textbf{RUNNING} maintains a
ring buffer of the latest \emph{n} (a moving window);
\textbf{GROWING} accumulates up to a maximum (the visible Law of
Large Numbers). The live visualization is the buffer's histogram
with a vertical mean line and a shaded confidence-interval band
whose width scales as $1 / \sqrt{n}$. Polyphonic input is supported
so a Methods pipeline scales to multiple parallel pseudo-experiments
without duplicating modules.

\begin{figure}[htbp]
\centering
\includegraphics[width=\linewidth]{figures/screenshots/anchor_sample_frame.png}
\caption{\textbf{Sample and Frame in action.} Left: \texttt{Seed}
(217) reproducibly initializes the pipeline. Centre: \texttt{Sample}
configured as Normal, showing the empirical histogram (blue bars)
converging on the theoretical PDF (gold curve) after $n = 128$
draws. Right: \texttt{Frame} in GROWING mode, $n = 4096$ --- the
sample mean has converged to $-0.00$ and the standard error has
collapsed to $0.004$ ($\approx 1/\sqrt{4096}$), visible as the
narrow CI band on the histogram.}
\label{fig:anchor-sample-frame}
\end{figure}

\textbf{\texttt{Regress}} reads pairs ($X$, $Y$) on a clock, stores
them in a buffer, and fits the line $Y = \alpha + \beta X$ by
minimizing the sum of squared residuals. The live scatterplot draws
the data, the fitted line, and a shaded confidence band whose width
tightens in the middle of the $X$ range and widens at the extremes
--- the trumpet shape that classically conveys the heteroscedasticity
of the prediction interval. The slope $\beta$, intercept $\alpha$,
$R^2$, and current residual $Y - (\hat\alpha + \hat\beta X)$ are
emitted as CV signals, so a downstream \texttt{Test} or
\texttt{Boot} module can take a hypothesis test or a bootstrap CI
on $\beta$ directly without leaving the rack.

\begin{figure}[htbp]
\centering
\includegraphics[width=\linewidth]{figures/screenshots/anchor_regress.png}
\caption{\textbf{Regress with a null-relationship pairing.} A
\texttt{Sample} (Uniform) feeds $X$; an independent \texttt{Sample}
(Normal) feeds $Y$. The resulting scatter has no genuine linear
relationship --- $\hat\beta \approx 0$ --- and the fitted line is
nearly horizontal. The shaded confidence band shows the trumpet
shape characteristic of an OLS prediction interval. Switching one
\texttt{Sample}'s preset introduces a genuine relationship and the
band straightens around a sloped line; the contrast lets a student
see what the band is actually showing.}
\label{fig:anchor-regress}
\end{figure}

\textbf{\texttt{Test}} computes a two-tailed \emph{t}-statistic
against H$_0$: $\mu = \mu_0$ (one-sample) or
H$_0$: $\mu_1 - \mu_2 = \delta_0$ (two-sample, Welch's \emph{t} for
unequal variances). The two-tailed \emph{p}-value is evaluated
\emph{exactly} through the regularized incomplete beta function via
Lentz's continued-fraction expansion \citep{press2007} --- the same
algorithm underlying R's \texttt{pt()} and SciPy's
\texttt{scipy.stats.t.sf()}, yielding machine-precision agreement at
any df $\geq 1$. The live visualization draws the
\emph{t}-distribution null with shaded rejection regions at the
user-selected $\alpha$, marks the observed \emph{t} as a vertical
line, and shows a REJECT / n.s.\ indicator; \emph{t}, \emph{p}, the
reject gate, and Cohen's \emph{d} effect size are emitted as CV
signals.

\begin{figure}[htbp]
\centering
\includegraphics[width=\linewidth]{figures/screenshots/anchor_test.png}
\caption{\textbf{Two-sample Welch \emph{t}-test with reproducible
seed.} \texttt{Seed} (544) fires two \texttt{Sample} modules; one
drives \texttt{Frame} for descriptives and both feed \texttt{Test}'s
SIG and SIG$_2$ inputs. The Test panel shows the null
\emph{t}-distribution with shaded rejection regions at the
user-selected $\alpha$, the observed $t = -3.01$ marked as a
vertical line in the left tail, and the REJECT indicator with the
supporting $p = 0.001$ readout. Because the entire pipeline is
seeded, the same patch opened on any machine reproduces this
\emph{t}, this \emph{p}, and this decision exactly.}
\label{fig:anchor-test}
\end{figure}

\textbf{\texttt{Boot}} is the simulation-based-inference workhorse.
Given a buffer of samples it draws \emph{B} resamples with
replacement and computes a user-selectable statistic (mean, median,
SD, variance) on each, so the \emph{empirical bootstrap
distribution} of the estimator becomes visible on the panel ---
\emph{this is} the sampling distribution made explicit, without
parametric assumption. The bias-corrected and accelerated (BCa)
interval \citep{efron1987} is reported alongside the point estimate
and the bootstrap SE; the panel surfaces the bias-correction $z_0$
and the acceleration $\hat a$ as live readouts, so a student can
watch the BCa correction shift the interval away from the
percentile baseline when the bootstrap distribution is skewed.
Patching a clock into the RSAMP input fires a fresh bootstrap on
every tick --- useful for visualizing the bootstrap's own
variability against a frozen sample.

The remaining ten Methods modules support the anchor pipeline:
\texttt{Lag} computes the autocorrelation function with Bartlett
bands; \texttt{Code} discretizes continuous CV into ordinal Likert
categories; \texttt{Tab} constructs contingency tables with the
exact $\chi^2$ \emph{p}-value via the regularized lower incomplete
gamma; \texttt{Strata} implements STL decomposition
\citep{cleveland1990}; \texttt{Cohort} performs online \emph{k}-means
quantization; \texttt{Factor} performs online PCA via Sanger's rule
\citep{sanger1989}; \texttt{Seed} is the reproducibility primitive;
\texttt{Tape} records and replays CV streams and supports CSV
export; \texttt{Gauge} and \texttt{Quantity} are a symmetric pair
that interpret CV voltages as real-world quantities and vice versa.
Complete per-parameter, per-input, per-output reference for all
fifteen modules is in the plugin manual.

The broader Empiria suite --- four additional plugins covering
agent-based social-science simulators (\textbf{Polis}), network
epidemiology (\textbf{Epi}), spatial-emergence models
(\textbf{Space}), and behavioral-economics models
(\textbf{Decisions}) --- shares the Methods panel grammar, the
common patch-cable conventions, and the same deterministic-seed
reproducibility story. We retain one cross-plugin Schelling
demonstration in Workflow A below to illustrate the composition;
substantive treatment of the broader suite is the subject of a
separate companion paper.

\section{Three worked examples}\label{three-worked-examples}

The three worked examples below illustrate the platform across its
intended audience range. Workflow A is a \emph{visceral}
demonstration suitable for an audience that has never met a
probability distribution; Workflow B is a \emph{precise}, seeded,
exportable analysis at a level a graduate methods student would
scrutinize; Workflow C is a \emph{verification} exercise
demonstrating that Empiria's empirical estimate of a closed-form
probability agrees with the analytic value. All three patches are
short (3--5 modules) and ship under \texttt{patches/} as
\texttt{.vcv} files; all three reproduce byte-identically on any
machine.

\subsection{Workflow A --- Watching segregation emerge}\label{workflow-a}

\begin{figure}[htbp]
\centering
\includegraphics[width=0.7\linewidth]{figures/patches/schelling.png}
\caption{\textbf{Starter patch A --- Schelling cross-plugin
demonstration.} A \texttt{Seed} primitive triggers the
\texttt{Schelling} reset input; \texttt{Frame} records the running
segregation index. Schelling is from the Space plugin (out of scope
for this paper) but is included here as a one-module illustration of
how \texttt{Frame} attaches to a substantive simulator rather than
to a parametric \texttt{Sample}.}
\label{fig:patch-schelling}
\end{figure}

A single Schelling module. Set the tolerance knob $\theta = 0.30$,
the occupancy knob to 0.85, balance to 0.50, and press SHUFFLE. The
grid starts as colour-noise --- red and blue cells randomly
intermixed. Over the next handful of ticks the boundaries
crystallize into large monochromatic clusters; the on-panel
segregation index rises from roughly 0.4 to roughly 0.75, and the
unhappy-agent fraction collapses toward zero.

The instructor asks the audience to predict, before dialing, what
will happen at $\theta = 0.20$ (essentially no preference) versus
$\theta = 0.55$ (strong preference). Sweeping $\theta$ live reveals:
at $\theta = 0.20$ the grid stays roughly mixed; at $\theta = 0.55$
segregation is total. The canonical Schelling lesson --- large
macro-pattern from mild micro-preference --- is conveyed without a
single line of mathematics, in under five minutes, to an audience
that may never have heard of agent-based modeling. The visceral,
knob-driven character of the demonstration is what makes it usable
in a high-school assembly or a public-engagement event as readily
as in a research seminar.

\subsection{Workflow B --- Seeded, reproducible, exportable}\label{workflow-b}

\begin{figure}[htbp]
\centering
\includegraphics[width=0.8\linewidth]{figures/patches/seeded_t_test.png}
\caption{\textbf{Starter patch B --- seeded two-sample Welch
\emph{t}-test.} \texttt{Seed} deterministically resets two
independent \texttt{Sample} modules so the same patch produces a
byte-identical \emph{t} statistic and \emph{p} value on any
machine. \texttt{Sample 1} feeds both the \texttt{Frame} summary and
\texttt{Test}'s SIG input; \texttt{Sample 2} feeds \texttt{Test}'s
SIG2 input.}
\label{fig:patch-ttest}
\end{figure}

The same modules can be used at a level that satisfies a graduate
methods student. We run a small-sample one-sample \emph{t}-test
against H$_0 = 0$, with the explicit goal of comparing Empiria's
\emph{exact} Student-\emph{t} \emph{p}-value against the
normal-approximation \emph{p}-value that many introductory teaching
tools still emit.

\textbf{Setup.} Place a \texttt{Seed} module with VALUE = 42; the
same integer on another machine reproduces byte-identical results.
Patch Seed's TRIG output into Sample's RESET and Test's RESET so
both downstream modules re-initialize in lockstep. Place
\texttt{Sample} with DIST = Normal, P1 ($\mu$) = 0.5, P2 ($\sigma$)
= 1.0 --- a deliberately ambiguous case for an $n = 8$ sample, where
the true mean is positive but small relative to $\sigma$. Place
\texttt{Test} with MODE = SNAPSHOT, $N = 8$, ALPHA = 0.05, H$_0 =
0$. Patch \texttt{Sample.SAMPLE} $\rightarrow$ \texttt{Test.SIG} and
press Seed's RANDOMIZE.

\textbf{Observable behavior.} Test's panel shows the observed
\emph{t} and the exact two-tailed \emph{p}-value. The null
\emph{t}-PDF visualization uses the correct Student's-\emph{t}
distribution at df = 7 --- visibly fatter-tailed than the standard
normal --- and the rejection regions are shaded according to the
chosen $\alpha$. The pedagogical claim is graduate-level: with df
= 7, the exact tail probability is materially larger than what a
normal-approximation tool would report, and the discrepancy can
flip a conclusion at $\alpha = 0.05$. The student can compare
Empiria's \emph{p} against R's
\texttt{pt(t, df = 7, lower.tail = FALSE) * 2} for any (\emph{t},
df) and confirm agreement to machine precision.

\textbf{External round-trip.} To carry the analysis outside the
platform, place a \texttt{Tape} module, set MODE = REC, LENGTH = 8,
patch \texttt{Sample.SAMPLE} $\rightarrow$ \texttt{Tape.SIG}, let it
fill, then right-click Tape $\rightarrow$ \textbf{Export buffer to
CSV\ldots}. The resulting CSV is a single column of the eight raw
draws, immediately loadable in R as
\texttt{d <- read.csv("sample.csv")\$ch1}; the standard
\texttt{t.test(d, mu = 0)} matches Empiria's \emph{t} and \emph{p}
to six decimal places. The same CSV, generated under the same Seed
value on a different machine, is byte-identical. This is the
workflow that makes Empiria-based assignments \emph{gradable}: the
instructor distributes a \texttt{.vcv} patch plus seed value,
students run it, export the CSV, and submit both their derived
statistic and the underlying data --- all of which the instructor
can verify reproduces.

Workflows A and B are the same tool in two modes: the rack stays
the same, the audience changes, the depth at which the underlying
mathematics is interrogated changes.

\subsection{Workflow C --- Verifiable computation against a closed-form}\label{workflow-c}

\begin{figure}[htbp]
\centering
\includegraphics[width=0.85\linewidth]{figures/patches/coin_flip_lln.png}
\caption{\textbf{Starter patch C --- coin-flip Law of Large
Numbers.} A near-Bernoulli \texttt{Sample} (Beta with
$\alpha = \beta \approx 0.05$) feeds \texttt{Frame} in GROWING mode
and \texttt{Boot}. \texttt{Gauge} with the Probability preset
reinterprets the running mean as a 0--1 quantity for legibility.}
\label{fig:patch-coin}
\end{figure}

Workflows A and B exhibit Empiria's pedagogical surface; Workflow C
is the \emph{rigor proof}. The simulation-based-inference
literature recommends that learners construct empirical sampling
distributions and check them against closed-form values where both
are available \citep{cobb2007, tintle2014}; doing so is also the
standard way to validate that a computational instrument behaves as
advertised. If Empiria's empirical estimate reproduces a textbook
combinatorial probability to within Monte Carlo error,
deterministically, on any operating system, then its claim to
methodological seriousness is concrete rather than rhetorical.

\textbf{The closed-form target.} For a fair coin, the probability
of four heads in a row in four flips is
$(1/2)^4 = 1/16 = 0.0625$ exactly. We estimate it empirically by
running 10\,000 independent four-flip trials.

\textbf{Setup.} \texttt{Seed} VALUE = 100; patch TRIG into Sample's
RESET and Tape's RESET. \texttt{Sample} DIST = Uniform, P1 = 0, P2
= 1. \texttt{Code} $K = 2$, LOW = 0, HIGH = 1 --- encodes each
uniform draw as 1 V (tails) or 2 V (heads) at threshold 0.5.
\texttt{Tape} LENGTH = 4096, MODE = REC, patched
\texttt{Code.CAT} $\rightarrow$ \texttt{Tape.SIG}. After Tape has
filled to 40\,000 samples (at the internal 200 Hz clock,
$\sim$\,200 seconds), switch MODE to PLAY and right-click Tape
$\rightarrow$ \textbf{Export buffer to CSV\ldots}.

\textbf{Verification in R.}

\begin{verbatim}
d     <- read.csv("coin_flips_seed100_n40000.csv")$ch1
flip  <- as.integer(d > 1.5)        # 1 = heads, 0 = tails
trial <- matrix(flip, nrow = 4)     # 4 flips / trial -> 10000 trials
phat  <- mean(colSums(trial) == 4)  # empirical fraction of HHHH
phat
#> [1] 0.0626

binom.test(sum(colSums(trial) == 4), ncol(trial), p = 1/16)
#> Exact binomial test ... p-value = 0.96  (cannot reject p = 1/16)
\end{verbatim}

The Monte Carlo standard error at 10\,000 trials with the true
$p = 0.0625$ is $\sqrt{p(1-p)/n} \approx 0.0024$. The observed
\texttt{phat} is within 0.5 SE of the true value; an exact binomial
test does not reject the null that the true probability is 1/16. A
collaborator running the same Empiria patch with Seed = 100 on a
different machine produces a byte-identical CSV, and therefore the
exact same \texttt{phat} to six decimal places.

The pedagogical lesson --- that closed-form combinatorial
probabilities have empirically verifiable counterparts ---
generalizes to harder targets (the probability of at least one run
of four heads in twenty flips, $\approx 0.184$, is reached by the
same patch with trial size changed to 20), and the modular
synthesizer turns out to be a perfectly adequate substrate for the
verification.

\section{Reproducibility, classroom integration, and evaluation}\label{repro-classroom}

Empiria treats classroom-scale reproducibility as a first-order
design property rather than as a downstream concern. Every random
module wraps \texttt{std::mt19937} (Mersenne Twister, period
$2^{19937}-1$) seeded from a per-module 32-bit seed preserved in
the patch JSON and broadcast on demand by the \texttt{Seed} module
\citep{matsumoto1998}. Patches are stored as human-readable JSON;
CSV exports from \texttt{Tape} are plain ASCII with six decimal
digits per sample. Three students on three different operating
systems running the same patch with the same Seed value produce
three byte-identical CSV files. Every analytical output that
Empiria produces has an external reference against which it can be
checked: R's \texttt{t.test()}, \texttt{lm()},
\texttt{boot::boot.ci()}, \texttt{pchisq()}; Python's
\texttt{scipy.stats}. Agreement is verified at the design stage and
is re-checkable by any user willing to perform the round-trip ---
exactly what Workflow C demonstrates.

The platform is usable across three classroom levels. In a
\textbf{high-school or general-education setting}, the platform
stands on its own as a hands-on demonstration tool --- no R, no
Python, no command line. A teacher can run a 45-minute session
built around a single patch (e.g.\ \texttt{Sample} $\rightarrow$
\texttt{Frame}, sweeping the \emph{N} knob to make the standard
error shrink visibly as $1/\sqrt{n}$), and students engage by
turning knobs and predicting what the trace will do. In an
\textbf{undergraduate introductory statistics or research-methods
course}, Empiria is positioned as a companion tool that visualizes
what R or Python is computing, not as a replacement. Students
install VCV Rack (free) and the Empiria \texttt{.vcvplugin} files
(total $<$ 1 MB) on personal laptops with no administrator
privileges required. In a \textbf{graduate research-methods
seminar}, the same modules become a fast-prototyping sandbox for
interrogating edge cases --- what does the null
\emph{t}-distribution look like at df = 3? how does the BCa
interval shift away from the percentile interval on a skewed
bootstrap? The same modules, knobs, and outputs serve all three
audiences; what differs is the depth at which the underlying
mathematics is interrogated.

A natural next step is empirical evaluation against a control
condition. We sketch a within-instructor, between-section randomized
comparison over two academic terms: parallel sections of the same
course, taught by the same instructor, randomly assigned to an
Empiria-augmented condition (Empiria plus the usual R workflow) or
a control (R workflow only, with comparable static figures and
Shiny applets to control for time on visualizations). Outcomes
would be measured by the \emph{Comprehensive Assessment of Outcomes
in a First Statistics course} (CAOS), with the \emph{Sampling
Variability}, \emph{Statistical Inference}, and \emph{Bivariate
Data} subscales as primary endpoints; a paper-based transfer task
scored blind by two graders; and a brief end-of-term confidence and
motivation survey. The most plausible failure mode of the platform
is \emph{activation cost} --- students who have never opened a
modular synthesizer may need 30 minutes of orientation before they
can drive the platform productively --- which the deployment plan
addresses with a dedicated orientation session and four worked
reference patches. A secondary risk is \emph{content shift}, which
the recommended additive (not substitutive) integration with the R
workflow is intended to control.

\section{Conclusion}\label{conclusion}

Empiria places fifteen statistical primitives into the
patch-cable grammar of a modular synthesizer so that the
inferential workflow --- from data-generating process through
estimation, hypothesis testing, and bootstrap confidence intervals
--- becomes a manipulable, visible, left-to-right signal chain on
the screen. The platform is intended to complement, not replace,
the analytical R/Python workflow that anchors current statistics
teaching: a student who has wired a \emph{t}-test together and
watched its null distribution shade in response to dialled-in data
carries a richer mental model into the eventual \texttt{t.test()}
call. Three properties --- seeded reproducibility, exact numerical
recipes for closed-form statistics, and portable JSON patches with
no external library dependencies --- make the platform defensible
as a \emph{reproducible analog computer} alongside its pedagogical
role.

The platform's most significant limitations are the absence of a
controlled classroom evaluation (the next step planned by the
author), the activation cost for users new to the modular interface,
and the polyphonic-cable cap of sixteen channels which limits agent
counts in the broader Empiria suite's social-science simulators.
None of these limitations is intrinsic to the design; each is a
planned-but-deferred next step.

\section{Acknowledgements}\label{acknowledgements}

The author thanks the VCV Rack developer community for the modular
plugin SDK and example modules, and acknowledges the long
open-source tradition of computational social science on which
Empiria draws.

\section{Disclosure statement}\label{disclosure-statement}

The author reports there are no competing interests to declare.
Empiria is released under the GPL-3.0 license through the author's
personal open-source imprint. The imprint is not a commercial
entity and yields no income to the author. No external funding
supported this work.

\section{Data availability statement}\label{data-availability-statement}

This paper describes a software platform rather than the analysis
of a specific dataset; the materials that support its results are
the source code, panel assets, reference patches, and documentation
that together constitute the platform. All such materials are
available under the GPL-3.0 license at
\url{https://github.com/kevisc/empiria}. The repository tag
\texttt{v2.0.0} corresponds to the version described in this paper
and is permanently archived at
\url{https://doi.org/10.5281/zenodo.20283281}. The repository
contains the C++17 source for all fifteen Methods modules, panel
SVG assets, the three starter patches as portable JSON
\texttt{.vcv} files, the patch-generator script
\texttt{tools/empiria\_patch\_gen.py}, the Methods Plugin Manual at
\texttt{docs/methods\_manual.md}, and the source for this paper.
